\documentclass[../../main.tex]{subfiles}
\begin{document}

{\bf [解答]} まず前半部について，$n\ge2$のとき，
\begin{align}
  f'_n(x)&=1+\dfrac{x}{1!}+\dfrac{x^2}{2!}+\dots+\dfrac{x^{n-1}}{(n-1)!}\nonumber \\
  &=f_{n-1}(x) \label{eq:1}
\end{align}
であるから示された．$\Box$

そこで後半部を考える．
\begin{center}
  P:方程式$f_n(x)=0$は，$n$が奇数ならばただ$1$つの実根をもつ．
  また$n$が偶数ならば$y=f_n(x)>0$である．
\end{center}
として，命題$P$を帰納的に示す．
まず$n=1,2$のとき，
\begin{align}
  \left\{
    \begin{array}{l}
      f_1(x)=1+x \\
      f_2(x)=1+x+\dfrac{x^2}{2}=\dfrac{1}{2}(x+1)^2+\dfrac{1}{2}
    \end{array}
    \right.
  \end{align}
  が$P$を満たすことは明らかである．そこで$k\in\mathbb{N}$に対し$n=2k-1,2k$での$P$の成立を仮定する．

  \subparagraph{(i)$n=2k+1$のとき}

  \cref{eq:1}および仮定から$f'_{2k+1}=f_{2k}(x)>0$だから$y=f_{2k+1}$は単調増加である．
  また$f_{2k+1}$は連続で$\lim_{x\to\pm\infty}f_{2k+1}(x)=\pm\infty$(複号同順)であるから，中間値の定理より
  $f_{2k+1}(x)=0$なる実数$x$がただひとつある．ここで$f_{2k+1}(0)=1$から$x\not=0$である．

  \subparagraph{(ii)$n=2(k+1)$のとき}

  \cref{eq:1}から$f'_{2(k+1)}(x)=f_{2k+1}(x)$であり，
  $n=2k+1$の時の考察から，$f_{2k+1}(a)=0$なる実数$a$がただひとつあって下表を得る．

  \begin{table}[htb]
    \centering
    \caption{$f_{2(k+1)}(x)$ の増減表}
    \label{tab:f_2k_plus_2}
    \begin{tabular}{|c|c|c|c|}
      \hline
      $x$ & $\cdots$ & $a$ & $\cdots$ \\ \hline
      $f'_{2(k+1)}(x)$ & $-$ & $0$ & $+$ \\ \hline
      $f_{2(k+1)}(x)$ & $\searrow$ & $f_{2(k+1)}(a)$ & $\nearrow$ \\ \hline
    \end{tabular}
  \end{table}
  故に
  \begin{align}
    f_{2(k+1)}(x)&\ge f_{2(k+1)}(a) \\
    &= f_{2k+1}(a)+\dfrac{a^{2(k+1)}}{2(k+1)!}   \\
    &=\dfrac{a^{2(k+1)}}{2(k+1)!}
    >0
  \end{align}
  である．
  \casesend

  以上から$n=2k+1,2(k+1)$でも$P$が成り立つことが示された．故に任意の$n$で$P$は成立する．
  $P$が示されたので，題意が従う．$\Box$

  \newpage
  \end{document}
